\begin{figure}[H]
    \centering
    \begin{tikzpicture}[x=2.45cm, y=1.25cm]

        \foreach \c in {0,1,2,3} {
            \node[blknpu, minimum width=1.30cm, minimum height=0.58cm] (r3\c) at (\c,2.10) {$(\c,3)$};
            \node[blknpu, minimum width=1.30cm, minimum height=0.58cm] (r2\c) at (\c,1.20) {$(\c,2)$};
            \node[blkdrop, minimum width=0.55cm, minimum height=0.44cm] (r1\c) at (\c,0.41) {\tiny \tl{mem}};
            \node[blknpu, minimum width=1.90cm, minimum height=0.48cm] (r0\c) at (\c,-0.34) {\tl{shim} $(\c,0)$};
        }

        \foreach \c in {0,1,2,3} {
            \draw[flowmute] ([xshift=-12pt]r0\c.north) -- ([xshift=-12pt]r2\c.south);
            \draw[flowmute] ([xshift=12pt]r2\c.south) -- ([xshift=12pt]r0\c.north);
            \draw[flowmute] ([xshift=-26pt]r0\c.north) -- ([xshift=-26pt]r3\c.south);
            \draw[flowmute] ([xshift=26pt]r3\c.south) -- ([xshift=26pt]r0\c.north);
        }

        \node[blkdrm, minimum width=9.4cm, minimum height=0.48cm] (ddr) at (1.5,-1.22)
            {κύρια μνήμη \tl{DDR5} (κοινή με τον \tl{host})};
        \foreach \c in {0,1,2,3} { \draw[flowmute] (r0\c.south) -- (r0\c.south|-ddr.north); }

        \draw[brace] (-0.30,2.50) -- (3.30,2.50);
        \node[notemute, font=\tiny, anchor=south] at (1.5,2.62)
            {8 υπολογιστικά πλακίδια: 4 στήλες $\times$ 2 πυρήνες};

        \node[notemute, anchor=west, align=left] at (3.42,1.65)
            {\tl{object\_fifo}\\ βάθος 2\\ (διπλή ενταμίευση)};
        \node[notemute, anchor=west, align=left] at (3.42,0.41) {δεν χρησιμοποιείται};
        \node[notemute, anchor=west, align=left] at (3.42,-0.34) {διεπαφή προς \tl{DDR}};

        \node[note, anchor=north, align=center] at (1.5,-1.62)
            {Ανά πυρήνα: 128 τεμάχια $\times$ 1024 στοιχεία \tl{int32}.
             Ανά παρτίδα: $8 \times 128 \times 1024 = 1\,048\,576$ στοιχεία.};

    \end{tikzpicture}
    \caption{Η τοπολογία πλακιδίων της υλοποίησης και οι ουρές \tl{object\_fifo} που
    τα συνδέουν. Κάθε στήλη διαθέτει ένα πλακίδιο διεπαφής (\tl{shim}) που
    επικοινωνεί με την κύρια μνήμη και δύο υπολογιστικά πλακίδια, καθένα από τα
    οποία συνδέεται με το \tl{shim} της στήλης του μέσω ενός ζεύγους ουρών βάθους
    δύο: μία για την είσοδο και μία για την έξοδο. Το βάθος δύο επιτρέπει στον
    πυρήνα να επεξεργάζεται ένα τεμάχιο ενώ το επόμενο μεταφέρεται. Τα πλακίδια
    μνήμης της ενδιάμεσης σειράς δεν χρησιμοποιούνται, καθώς κάθε πυρήνας
    καταναλώνει ξένο προς τους υπόλοιπους τμήμα της εισόδου και δεν απαιτείται
    διαμοιρασμός ή αναδιανομή δεδομένων μεταξύ τους.}
    \label{fig:npu-topology}
\end{figure}
